\clearpage\pagestyle{otMain}\setlength{\otBlocEspacement}{2.5mm}\fontsize{11}{15}\selectfont\corrigesChapitre{14}{Thalès : triangles emboîtés}
\begin{blocLivre}{cor-14-1}\begin{corrige}{1}{Un rapport simple}\demarche{Relier petit et grand côtés correspondants.}
Les points sont alignés dans l’ordre et les droites parallèles. Thalès donne $AN/10=AM/AB=4/8=1/2$. Donc $\res{AN=5\ \mathrm{cm}}$.
\label{fig-cor-14-1}\figureGeo[corrige]{th-01}
\end{corrige}\end{blocLivre}
\begin{blocLivre}{cor-14-2}\begin{corrige}{2}{Le segment parallèle}\demarche{Utiliser le rapport de réduction.}
Par Thalès, $MN/9=8/12=2/3$. Donc $MN=9\times2/3=\res6$ cm. La longueur est plus petite que BC, ce qui est cohérent avec les triangles emboîtés.
\label{fig-cor-14-2}\figureGeo[corrige]{th-02}
\end{corrige}\end{blocLivre}
\begin{blocLivre}{cor-14-3}\begin{corrige}{3}{Retrouver le grand côté}\demarche{Inverser le passage du grand vers le petit.}
$MN/BC=AM/AB=1/3$. Ainsi $4/BC=1/3$, donc $\res{BC=12\ \mathrm{cm}}$. On multiplie ici le petit côté par 3 pour revenir au grand.
\label{fig-cor-14-3}\figureGeo[corrige]{th-03}
\end{corrige}\end{blocLivre}
\begin{blocLivre}{cor-14-4}\begin{corrige}{4}{Une réciproque}\demarche{Vérifier les deux conditions d’alignement et les rapports.}
$AM/AB=6/10=3/5$ et $AN/AC=4{,}8/8=3/5$. Les points sont sur les segments correspondants et les rapports sont égaux. La réciproque de Thalès donne $\res{(MN)\parallel(BC)}$.
\label{fig-cor-14-4}\figureGeo[corrige]{th-04}
\end{corrige}\end{blocLivre}
\begin{blocLivre}{cor-14-5}\begin{corrige}{5}{Des rapports à simplifier}\demarche{Comparer exactement deux fractions.}
$6/9=2/3$ et $8/12=2/3$. Avec les points sur les segments, la réciproque s’applique : les droites sont parallèles.
\label{fig-cor-14-5}\figureGeo[corrige]{th-05}
\end{corrige}\end{blocLivre}
\begin{blocLivre}{cor-14-6}\begin{corrige}{6}{Un parallélisme impossible}\demarche{Utiliser le théorème pour exclure une hypothèse.}
$AM/AB=3/5$ et $AN/AC=8/15$. Or $3/5=9/15\ne8/15$. Si les droites étaient parallèles, Thalès imposerait l’égalité. Elles ne sont donc pas parallèles.
\label{fig-cor-14-6}\figureGeo[corrige]{th-06}
\end{corrige}\end{blocLivre}
\begin{blocLivre}{cor-14-7}\begin{corrige}{7}{Calculer le bon segment}\demarche{Reconstituer la longueur depuis le sommet commun.}
$AM=AB-BM=12-3=9$ cm. Thalès donne $AN/8=9/12=3/4$, donc $\res{AN=6\ \mathrm{cm}}$. Utiliser BM à la place de AM donnerait le mauvais rapport.
\label{fig-cor-14-7}\figureGeo[corrige]{th-07}
\end{corrige}\end{blocLivre}
\begin{blocLivre}{cor-14-8}\begin{corrige}{8}{Un coefficient décimal}\demarche{Utiliser une seule fois le coefficient commun.}
Le rapport petit sur grand vaut $4/10=0{,}4$. Donc $AN=0{,}4\times20=\res8$ cm et $MN=0{,}4\times15=\res6$ cm. Le même coefficient s’applique aux côtés correspondants.
\label{fig-cor-14-8}\figureGeo[corrige]{th-08}
\end{corrige}\end{blocLivre}
\begin{blocLivre}{cor-14-9}\begin{corrige}{9}{Un protocole de construction}\demarche{Construire le parallélisme avant de calculer.}
Construire ABC par CAC, placer M sur $[AB]$, puis tracer la parallèle à BC passant par M. Son intersection avec AC est N. Le rapport est $3/9=1/3$ ; Thalès donne $\res{AN=4\ \mathrm{cm}}$.
\label{fig-cor-14-9}\figureGeo[corrige]{th-09}
\end{corrige}\end{blocLivre}
\begin{blocLivre}{cor-14-10}\begin{corrige}{10}{Une proportion inverse}\demarche{Isoler le petit côté dans le rapport.}
$AM/7=MN/BC=5/14$. Donc $AM=7\times5/14=\res{2{,}5\ \mathrm{cm}}$. La valeur est inférieure à AB, comme l’exige M sur le segment.
\label{fig-cor-14-10}\figureGeo[corrige]{th-10}
\end{corrige}\end{blocLivre}
\begin{blocLivre}{cor-14-11}\begin{corrige}{11}{Une hauteur par proportion}\demarche{Justifier le parallélisme par les verticales.}
Les verticales MN et BC sont parallèles, et les deux alignements sont donnés. Par Thalès, $MN/BC=AM/AB=1/5$. Donc $BC=1{,}5\times5=\res{7{,}5\ \mathrm m}$. Le modèle suppose les alignements annoncés.
\label{fig-cor-14-11}\figureGeo[corrige]{th-11}
\end{corrige}\end{blocLivre}
\begin{blocLivre}{cor-14-12}\begin{corrige}{12}{Un morceau restant}\demarche{Calculer depuis A puis soustraire au grand côté.}
$AN/10=3/8$, d’où $AN=30/8=\res{3{,}75}$ cm. Comme N est entre A et C, $NC=AC-AN=10-3{,}75=\res{6{,}25}$ cm.
\label{fig-cor-14-12}\figureGeo[corrige]{th-12}
\end{corrige}\end{blocLivre}
\begin{blocLivre}{cor-14-13}\begin{corrige}{13}{Justifier puis calculer}\demarche{Séparer la réciproque du théorème direct.}
$4/10=6/15=2/5$. La réciproque donne $(MN)\parallel(BC)$. On peut alors utiliser le théorème direct : $MN/12=2/5$, donc $\res{MN=4{,}8\ \mathrm{cm}}$. Le parallélisme est prouvé avant le calcul.
\label{fig-cor-14-13}\figureGeo[corrige]{th-13}
\end{corrige}\end{blocLivre}
\begin{blocLivre}{cor-14-14}\begin{corrige}{14}{Partager un segment}\demarche{Transférer un partage régulier par des parallèles.}
$AM/AC=3/5$. Les triangles AMN et ACB sont emboîtés et MN est parallèle à CB. Thalès donne $AN/AB=3/5$, donc $\res{AN=6\ \mathrm{cm}}$. La méthode construit les trois cinquièmes d’un segment sans mesurer directement cette fraction.
\label{fig-cor-14-14}\figureGeo[corrige]{th-14}
\end{corrige}\end{blocLivre}
\begin{blocLivre}{cor-14-15}\begin{corrige}{15}{Deux réductions emboîtées}\demarche{Soustraire les aires de deux triangles emboîtés.}
Thalès donne $AN=16\times3/12=4$ cm et $AQ=16\times9/12=12$ cm. Les triangles AMN et APQ sont rectangles en A. Leurs aires valent $3\times4/2=6$ et $9\times12/2=54$ cm$^2$. La bande M-N-Q-P a pour aire $54-6=\res{48\ \mathrm{cm}^2}$.
\label{fig-cor-14-15}\figureGeo[corrige]{th-15}
\end{corrige}\end{blocLivre}
